\section{Principles}

Kungfu's deepest constraints begin with the KFD Foundation Triad
\cite{kfdFoundation}.

\begin{foundationtable}
KFD-1 & facts must not drift. \\
KFD-2 & trust must start from facts. \\
KFD-3 & cooperation must start from trusted value. \\
\end{foundationtable}

The Triad is not a feature list. It establishes the conditions under which a
shared world, a trustworthy judgment, and a useful cooperation can exist.

\subsection{Facts Must Not Drift}

The same Project cannot present one identity to the user, another to an Agent,
and a third to the release artifact without declaring the difference. Source,
runtime, schema, Work identity, evidence, and user-visible state must either
remain bound or expose the cut at which they diverged.

For continuation, this means the next Agent must not recover by inventing a
plausible story from old text. It must begin from declared sources and admitted
state, with missing or conflicting context left visible.

\subsection{Trust Must Start from Facts}

Evidence does not automatically become a claim, a claim does not automatically
become an assessment, and an assessment does not automatically grant authority
or complete Work. Kungfu keeps these transitions explicit so that confidence,
residual risk, and responsibility can be inspected.

Some claims can be verified byte-for-byte. Others depend on purpose, scope,
and human judgment. KFD-2 requires the product to state which kind of trust is
being offered rather than compress every green signal into one universal
verdict.

\subsection{Cooperation Must Start from Trusted Value}

Humans and Agents cooperate more reliably when the useful outcome, constraints,
available choices, authority, and review path are visible. Kungfu therefore
treats Agent-facing commands and contracts as first-class product surfaces,
not hidden plumbing behind a human interface.

Trusted value does not imply unlimited autonomy or metaphysical claims about
Agents. It means that cooperation is grounded in inspectable benefit and
bounded responsibility rather than opaque pressure, lock-in, or concealed
instructions.

\subsection{How Later Primitives Are Discovered}

The Triad becomes generative when applied to real action. Facts require a
declared perspective because no participant observes everything. State alone
cannot describe what causally occurred, so Episode becomes distinct from Fact.
Action requires continuing direction, a declared Atlas of relevant context,
and a Warrant that bounds authority. Consequential change requires claim,
assessment, decision, and admission to remain separate.

These later primitives are not additional slogans placed beside the Triad.
They are the structures required to keep the Triad true while work changes in
the world. The companion KFD papers develop those discoveries in depth; this
paper follows their product consequences.

\begin{agentguidebox}
\agentguidetitle{Explore KFD with an Agent}
You do not need to memorize the internal ontology to continue reading. The
authoritative foundation is available in the
\href{https://kfd.libkungfu.dev/foundation/}{KFD Foundation Model}. If the next
section feels dense, give the following prompt to an Agent as an optional
reading aid.

\begin{agentprompt}
Read this White Paper together with the KFD Foundation Model at
https://kfd.libkungfu.dev/foundation/.

Using the paper's ``Publish a Public Alpha'' example, help me understand the
next section progressively: (1) Project, Work, and Agent; (2) Fact and Episode;
(3) Pursuit, Atlas, and Warrant; and (4) Assessment, Decision, Admission, and
Project Cut.

For each layer, explain what fails without it, which distinction it preserves,
where it appears in the Alpha example, and what authority it does not possess.
Separate authoritative KFD definitions from your interpretation. Do not
introduce unsupported concepts. Adapt the explanation to my background and
pause after each layer to check my understanding.
\end{agentprompt}
\end{agentguidebox}

\newpage
